% Generated draft: requires mathematical and source review.
\documentclass[11pt]{article}
\usepackage[margin=25mm]{geometry}
\usepackage{fontspec}
\setmainfont{Latin Modern Roman}
\usepackage{amsmath,amssymb,mathtools,unicode-math}
\setmathfont{Latin Modern Math}
\usepackage{enumitem,needspace,xurl,hyperref}
\setlength{\parindent}{0pt}
\setlength{\parskip}{5pt}
\setlength{\emergencystretch}{3em}
\newcommand{\N}{\mathbb{N}}
\newcommand{\Z}{\mathbb{Z}}
\newcommand{\Q}{\mathbb{Q}}
\newcommand{\R}{\mathbb{R}}
\newcommand{\C}{\mathbb{C}}
\newcommand{\F}{\mathbb{F}}
\newcommand{\A}{\mathbb{A}}
\newcommand{\PP}{\mathbb{P}}
\newcommand{\E}{\mathbb{E}}
\newcommand{\eps}{\varepsilon}
\newcommand{\dd}{\,\mathrm{d}}
\begin{document}
\section*{UnsolvedMath ID: 800966}
\textbf{Source problem code:} catalogue entry
\subsection{Definitions and mathematical statement}
The supplied record refers to two definitions, labelled (a) and (b), which are absent from the record. Write \(I_{(b)}(M,g)\) provisionally for the invariant defined in the missing clause (b), and write \(\mathcal S\) for the intended class of surfaces with metrics. Neither \(I_{(b)}\) nor \(\mathcal S\) can be defined precisely from the supplied material. The notation \(\operatorname{area}(M,g)\) denotes the area of \(M\) with respect to \(g\).

The question attributed in context to Guth asks whether there exists a constant \(C>0\), independent of the surface and its metric and, in particular, independent of its genus, such that
\[
\forall (M,g)\in\mathcal S,\qquad I_{(b)}(M,g)\leq C\sqrt{\operatorname{area}(M,g)}.
\]
This is only a schematic statement of the intended question: the missing definition (b) and the missing hypotheses specifying \(\mathcal S\) must be recovered before it constitutes a precise mathematical problem.

For context, the supplied source attributes to Balacheff and Sabourau, under the reference label [Bal-Sab2010], the bound
\[
\operatorname{dias}(M,g)\leq C\sqrt{\operatorname{genus}(M)}\sqrt{\operatorname{area}(M,g)}
\]
for the diastole defined in the missing clause (a), and reports that its dependence on genus is optimal. This contextual bound concerns the invariant in (a); the question concerns the invariant in (b). The definition of that diastole, the hypotheses of the bound, and the intended genus range are not supplied. In particular, the displayed right-hand side vanishes at genus zero under the usual genus convention, so its applicability in that case requires clarification.
\subsection{Short English statement}
The question attributed in context to Guth asks whether the invariant from the missing definition (b) has an upper bound by a universal constant times the square root of area, independently of genus. Its definition and the allowed surfaces and metrics are absent, so the supplied record does not determine a complete precise statement.
\subsection{Sources}
\begin{itemize}
\item Local UnsolvedMath registry, permanent ID \(800966\), source problem code \texttt{catalogue entry}; supplied title: \emph{Open Problems and Questions about Geodesics --- Question 47 (Guth)}. The registry labels the problem open; this is a supplied status, not an independently verified assessment.
\item Burns and Matveev, \emph{Open Problems and Questions about Geodesics}, source list dated 2019 in the record, Question 47: \url{https://arxiv.org/abs/1308.5417}. The registry lists Burns and Matveev as proposers and gives Guth as attribution/context. The supplied provenance records access on 2026-10-16 and extraction by \texttt{arxiv\_latex}; these provenance claims have not been independently verified here.
\item [Bal-Sab2010]: the supplied excerpt attributes the contextual diastolic inequality and optimal genus dependence to Balacheff and Sabourau. No full bibliographic entry or separate link for this reference is supplied.
\end{itemize}
\end{document}
